\documentclass[]{scrartcl}


%packages
\usepackage{amsmath}
\usepackage{xcolor}
\usepackage{hyperref}
\hypersetup{
	colorlinks   = true, %Colours links instead of ugly boxes
	urlcolor = teal
}

%opening
\title{Notes on Kalman filter for PTA analysis}
\author{T. Kimpson}

\begin{document}

\maketitle
\begin{abstract}
	This note defines the Kalman ``machinery" for the state-space formulation of PTA analysis in terms of timing residuals. It is heavily based on work by P. Meyers and the \textsc{minnow} package, with extensions to handle multiple pulsars and the influence of the stochastic GW background, with contributions from A. Vargas. This work builds on our previous papers, continuing from their state-space formulation in the frequency domain.
\end{abstract}

\section{Requirements}

For the purposes of running state-space algorithms, we define the following:

\begin{itemize}
	\item $\boldsymbol{X}$ : the state vector, dimension $n_X$
	\item $\boldsymbol{Y}$ : the observation vector, dimension $n_Y$
	\item $\boldsymbol{F}$ : the state transition matrix, dimension $n_X \times n_X$
	\item $\boldsymbol{H}$ : the measurement matrix, dimension $n_Y \times n_X$
	\item $\boldsymbol{Q}$ : the process noise covariance matrix, dimension $n_X \times n_X$
	\item $\boldsymbol{R}$ : the measurement noise covariance matrix, dimension $n_Y \times n_Y$
\end{itemize}

The goal of this note is to define the above matrices.



\section{System definition}

We have a set of $N$ pulsars. Each pulsar has $4 + M^{(n)}$ states, where $M^{(n)}$ is the number of timing model parameters for the $n$-th pulsar. The states evolve according to the following set of dynamical equations.
\begin{align}
	\frac{d}{dt} \delta \phi^{(n)} &= \delta f^{(n)}, \\
	\frac{d}{dt} \delta f^{(n)} &= -\gamma_p^{(n)} \delta f^{(n)} + \chi_p^{(n)}(t), \\
	\frac{d}{dt} r^{(n)} &= a^{(n)}, \\
	\frac{d}{dt} a^{(n)} &= -\gamma_a a^{(n)} + \chi_{\mathrm a}^{(n)}(t), \\
	\frac{d}{dt} \delta \epsilon^{(n)}_{m} &= \chi_{\epsilon_{m}}^{(n)}(t) \quad \forall m \in [1, M^{(n)}].
\end{align}
where $\chi_{i}$ denotes a white noise stochastic process and $\gamma_i$ are daming constants. The processes  $\chi_p^{(n)}$ and  $\chi_\epsilon^{(n)}$ are uncorrelated between different $n$. Their statistics are defined by

\begin{align}
	\langle \chi_p^{(n)}(t) \rangle &= 0 \ , \\
	\langle \chi_p^{(n)}(t) \chi_p^{(n')}(t') \rangle &= [\sigma_p^{(n)}]^2 \delta_{n,n'} \delta(t - t') \ .
\end{align}
The factor $\delta_{n,n'}$ in the right-hand side ensures that the non-GW-induced spin fluctuations in distinct pulsars $n\neq n'$ are uncorrelated, as expected physically. 

\begin{align}
	\langle \chi_{\epsilon_m}^{(n)}(t) \rangle &= 0 \ , \\
	\langle \chi_{\epsilon_m}^{(n)}(t) \chi_{\epsilon_m'}^{(n')}(t') \rangle &= [\sigma_{\epsilon_m}^{(n)}]^2 \delta_{n,n'} \delta(t - t') \ .
\end{align}

Note that I am not sure if it makes sense for the timing model parameter errors to be treated in this way; perhaps there is some correlation between the parameters for a single pulsar . For now we just assume that they are all uncorrelated for simplicity. \newline 

The correlation between pulsars manifests due to the GW, i.e. the $a^{(n)}$ state. The ensemble statistics of $\chi_{\mathrm a}^{(n)}$ are
\begin{align}
	\langle \chi^{(n)}_{\mathrm a}(t) \rangle &= 0 \ , 	\label{eq:xieqn1} \\
	\langle \chi^{(n)}_{\mathrm a}(t) \, \chi^{(n')}_{\mathrm a}(t') \rangle &= \left[\sigma^{(n,n')}_{\mathrm a}\right]^2 \delta(t - t') \ .	\label{eq:xieqn2}
\end{align}
with
\begin{eqnarray}
	\left[\sigma^{(n,n')}_{\mathrm a}\right]^2 = \frac{\langle h^2\rangle}{6} \gamma_{\mathrm a} \, \Gamma \left[ \theta^{(n,n')} \right] \, , \label{eq:sigma_a_expression}
\end{eqnarray}
where $\langle h^2 \rangle$ is the mean square GW strain, $\theta^{(n,n')}$ is the angle between the $n$-th and $n'$-th pulsars, and one has
\begin{eqnarray}
	\Gamma\left[\theta^{(n,n')} \right] =  \frac{3}{2} x_{n n'} \ln x_{n n'}  -\frac{x_{n n'} }{4}+\frac{1}{2} + \frac{1}{2} \delta_{n n'}\label{eq:correlation} \, ,
\end{eqnarray}
with $x_{nn'} = \left[1 - \cos \theta^{(n,n')}\right]/2$. \newline 

The states described by the SDEs are related to the observables, $\boldsymbol{Y} = \left(\delta t \right) $ via a measurement equation
\begin{equation}
	\delta t^{(n)} = \frac{\delta \phi^{(n)}}{f_0^{(n)}} + \mathbf{M}^{(n)} \boldsymbol{\delta \epsilon}^{(n)} - r^{(n)}
\end{equation}
where $\mathbf{M}^{(n)}$ is the design matrix of partial derivatives of the phases with respect to the timing model parameters. It is provided by standard pulsar timing libraries such as TEMPO and is a (row)vector of dimension $M^{(n)}$ (note the clumsy notation; $\mathbf{M}^{(n)}$ is a vector, the normal $M^{(n)}$ is a scalar which sets the dimension of $\mathbf{M}^{(n)}$) 
The quantity $\boldsymbol{\delta \epsilon}^{(n)}$ is a (column) vector, also of length $M^{(n)}$, which holds all the $\delta \epsilon^{(n)}_{m}$ terms for pulsar $n$.
The quantity $f_0^{(n)}$ is the nominal spin frequency of pulsar $n$.







\section{Block organisation}


The structure of the system invictes a block organisation of the state vector.

There are three blocks which are uncorrelated.
\begin{itemize}
	\item The spin-down states,$\delta \phi^{(n)}$ and $\delta f^{(n)}$.
	\item The GW states, $r^{(n)}$ and $a^{(n)}$.
	\item The timing model parameters, $\delta \epsilon^{(n)}_{m}$.
\end{itemize}

We will note these as $\boldsymbol{X}_{\text{spin}}$, $\boldsymbol{X}_{\text{gw}}$, and $\boldsymbol{X}_{\text{tm}}$, which stretch across all pulsars.
So for example $\boldsymbol{X}_{\text{spin}}  = \left(\delta \phi^{(1)}, \delta f^{(1)}, \delta \phi^{(2)}, \delta f^{(2)}, \cdots, \delta \phi^{(N)}, \delta f^{(N)}\right)$.


During the predict step of the Kalman filter, the block states can be evolved independently of each other
This is useful and allows us to work with small matrices. We will now derive the $F$ and $Q$ matrices for each block.



\subsection{GW block}
We consider a system of $N$ pulsars, each with state
\begin{equation}
x^{(n)} = \begin{pmatrix} r^{(n)} \\[1mm] a^{(n)} \end{pmatrix}, \quad n=1,\ldots,N,
\end{equation}
whose continuous--time dynamics are given by
\begin{equation}
\frac{d}{dt}r^{(n)} = a^{(n)}, \qquad \frac{d}{dt}a^{(n)} = -\gamma_a\,a^{(n)} + \chi_a^{(n)}(t).
\end{equation}
In state--space form this becomes
\begin{equation}
\frac{d}{dt} x^{(n)} = A\,x^{(n)} + B\,\chi_a^{(n)}(t),
\end{equation}
with
\begin{equation}
A=\begin{pmatrix} 0 & 1 \\[0.5mm] 0 & -\gamma_a \end{pmatrix}, \quad
B=\begin{pmatrix} 0 \\[0.5mm] 1 \end{pmatrix}.
\end{equation}

The exact discretisation over a time step $\Delta t$ is
\begin{equation}
x_{k+1}^{(n)} = e^{A\Delta t}\,x_k^{(n)} + \eta_k^{(n)},
\end{equation}
with
\begin{equation}
\eta_k^{(n)} = \int_0^{\Delta t} e^{A(\Delta t-s)}\,B\,\chi_a^{(n)}(t_k+s)\,ds.
\end{equation}
A straightforward calculation shows that the state--transition (or fundamental) matrix is
\begin{equation}
F_{\text{block}} \equiv e^{A\Delta t} =
\begin{pmatrix}
1 & \dfrac{1-e^{-\gamma_a\Delta t}}{\gamma_a} \\[1mm]
0 & e^{-\gamma_a\Delta t}
\end{pmatrix}.
\end{equation}

For the noise we define the $2\times2$ block
\begin{equation}
Q_{\text{block}} = \begin{pmatrix}
q_{11} & q_{12}\\[1mm]
q_{12} & q_{22}
\end{pmatrix},
\end{equation}
where
\begin{equation}
q_{22} = \frac{1-e^{-2\gamma_a\Delta t}}{2\gamma_a},
\end{equation}
\begin{equation}
q_{12} = \frac{1}{\gamma_a^2}\left[\left(1-e^{-\gamma_a\Delta t}\right)-\frac{1-e^{-2\gamma_a\Delta t}}{2}\right],
\end{equation}
\begin{equation}
q_{11} = \frac{1}{\gamma_a^3}\left[\Delta t-\frac{2}{\gamma_a}\left(1-e^{-\gamma_a\Delta t}\right)+\frac{1-e^{-2\gamma_a\Delta t}}{2\gamma_a}\right].
\end{equation}

Stacking the states of all $N$ pulsars into
\begin{equation}
X = \begin{pmatrix} x^{(1)} \\[1mm] x^{(2)} \\[1mm] \vdots \\[1mm] x^{(N)} \end{pmatrix},
\end{equation}
the full discrete system is
\begin{equation}
X_{k+1} = F\,X_k + \eta_k,
\end{equation}
with the state--transition matrix given by
\begin{equation}
F = I_N \otimes F_{\text{block}},
\end{equation}
where $I_N$ is the $N\times N$ identity matrix and $\otimes$ denotes the Kronecker product.

Similarly, if we define $\Sigma$ as the $N\times N$ matrix with entries
\begin{equation}
\Sigma_{n,n'} = \Bigl[\sigma^{(n,n')}_a\Bigr]^2,
\end{equation}
then the full noise covariance is written as
\begin{equation}
Q = \Sigma \otimes Q_{\text{block}}.
\end{equation}





\subsection{Spin-down block}


We consider a system of $N$ pulsars, each with state
\begin{equation}
x^{(n)} = \begin{pmatrix} \delta\phi^{(n)} \\[1mm] \delta f^{(n)} \end{pmatrix}, \quad n=1,\ldots,N.
\end{equation}
The continuous--time dynamics are given by
\begin{equation}
\frac{d}{dt}\delta\phi^{(n)} = \delta f^{(n)}, \quad \frac{d}{dt}\delta f^{(n)} = -\gamma_p^{(n)}\,\delta f^{(n)} + \chi_p^{(n)}(t).
\end{equation}
with noise statistics
\begin{equation}
\langle \chi_p^{(n)}(t) \rangle = 0,
\end{equation}
\begin{equation}
\langle \chi_p^{(n)}(t) \, \chi_p^{(n')}(t') \rangle = \left[\sigma_p^{(n)}\right]^2 \delta_{n,n'} \delta(t-t').
\end{equation}

For each pulsar the dynamics can be written in state--space form as
\begin{equation}
\frac{d}{dt}x^{(n)} = A^{(n)}\,x^{(n)} + B\,\chi_p^{(n)}(t),
\end{equation}
with
\begin{equation}
A^{(n)} = \begin{pmatrix} 0 & 1 \\[1mm] 0 & -\gamma_p^{(n)} \end{pmatrix}, \quad
B = \begin{pmatrix} 0 \\[1mm] 1 \end{pmatrix}.
\end{equation}

Stacking the states of all pulsars into
\begin{equation}
X = \begin{pmatrix} x^{(1)} \\[1mm] x^{(2)} \\[1mm] \vdots \\[1mm] x^{(N)} \end{pmatrix},
\end{equation}
the full continuous system is
\begin{equation}
\frac{d}{dt} X = \mathcal{A}\, X + \Xi(t),
\end{equation}
where
\begin{equation}
\mathcal{A} = \operatorname{diag}\Bigl(A^{(1)},\,A^{(2)},\,\ldots,\,A^{(N)}\Bigr),
\end{equation}
\begin{equation}
\Xi(t) = \begin{pmatrix} B\,\chi_p^{(1)}(t) \\[1mm] B\,\chi_p^{(2)}(t) \\[1mm] \vdots \\[1mm] B\,\chi_p^{(N)}(t) \end{pmatrix}.
\end{equation}

\bigskip
\textbf{Discretisation:}

For a time step $\Delta t$, the exact solution for the $n$th pulsar is
\begin{equation}
x_{k+1}^{(n)} = e^{A^{(n)}\Delta t}\,x_k^{(n)} + \eta_k^{(n)},
\end{equation}
with
\begin{equation}
\eta_k^{(n)} = \int_0^{\Delta t} e^{A^{(n)}(\Delta t-s)}\,B\,\chi_p^{(n)}(t_k+s)\,ds.
\end{equation}
A straightforward calculation shows that
\begin{equation}
F^{(n)} \equiv e^{A^{(n)}\Delta t} =
\begin{pmatrix}
1 & \dfrac{1-e^{-\gamma_p^{(n)}\Delta t}}{\gamma_p^{(n)}} \\[1mm]
0 & e^{-\gamma_p^{(n)}\Delta t}
\end{pmatrix}.
\end{equation}

Since each pulsar is independent, the full discrete state--transition matrix is
\begin{equation}
F = \operatorname{diag}\Bigl(F^{(1)},\,F^{(2)},\,\ldots,\,F^{(N)}\Bigr).
\end{equation}
Thus, the overall discrete system is
\begin{equation}
X_{k+1} = F\,X_k + \eta_k,
\end{equation}
with
\begin{equation}
\eta_k = \begin{pmatrix} \eta_k^{(1)} \\[1mm] \eta_k^{(2)} \\[1mm] \vdots \\[1mm] \eta_k^{(N)} \end{pmatrix}.
\end{equation}

For the noise, the covariance for the $n$th pulsar is given by
\begin{equation}
Q^{(n)} = \langle \eta_k^{(n)}\,\eta_k^{(n)\top} \rangle = \left[\sigma_p^{(n)}\right]^2 \int_0^{\Delta t} e^{A^{(n)}s}\,B\,B^\top\,e^{A^{(n)\top}s}\,ds.
\end{equation}
A direct evaluation yields
\begin{equation}
Q^{(n)} = \left[\sigma_p^{(n)}\right]^2 
\begin{pmatrix}
\displaystyle \frac{1}{\left[\gamma_p^{(n)}\right]^3}\left[\Delta t-\frac{2}{\gamma_p^{(n)}}
\left(1-e^{-\gamma_p^{(n)}\Delta t}\right)+\frac{1-e^{-2\gamma_p^{(n)}\Delta t}}{2\gamma_p^{(n)}}\right] & \displaystyle \frac{1}{\left[\gamma_p^{(n)}\right]^2}\left[\left(1-e^{-\gamma_p^{(n)}\Delta t}\right)-\frac{1-e^{-2\gamma_p^{(n)}\Delta t}}{2}\right] \\[2mm]
\displaystyle \frac{1}{\left[\gamma_p^{(n)}\right]^2}\left[\left(1-e^{-\gamma_p^{(n)}\Delta t}\right)-\frac{1-e^{-2\gamma_p^{(n)}\Delta t}}{2}\right] & \displaystyle \frac{1-e^{-2\gamma_p^{(n)}\Delta t}}{2\gamma_p^{(n)}}
\end{pmatrix}.
\end{equation}

Since the noise in different pulsars is uncorrelated, the full discrete noise covariance is
\begin{equation}
Q = \operatorname{diag}\Bigl(Q^{(1)},\,Q^{(2)},\,\ldots,\,Q^{(N)}\Bigr).
\end{equation}


\subsection{Timing model block}


We consider the state
\begin{equation}
X = \bigl(\delta\boldsymbol{\epsilon}^{(1)},\,\delta\boldsymbol{\epsilon}^{(2)},\,\dots,\,\delta\boldsymbol{\epsilon}^{(N)}\bigr),
\end{equation}
where for pulsar \(n\) the state is given by
\begin{equation}
\delta\boldsymbol{\epsilon}^{(n)} = \begin{pmatrix}
\delta\epsilon_1^{(n)} \\
\delta\epsilon_2^{(n)} \\
\vdots \\
\delta\epsilon_{M^{(n)}}^{(n)}
\end{pmatrix}.
\end{equation}
The continuous--time dynamics are defined by
\begin{equation}
\frac{d}{dt}\delta\epsilon_m^{(n)} = \chi_{\epsilon_m}^{(n)}(t) \quad \text{for } m=1,\ldots,M^{(n)}.
\end{equation}
Since there is no drift term (i.e. no dependence on \(\delta\boldsymbol{\epsilon}^{(n)}\) itself), we can write the dynamics for pulsar \(n\) as
\begin{equation}
\frac{d}{dt}\delta\boldsymbol{\epsilon}^{(n)} = I_{M^{(n)}}\,\boldsymbol{\chi}_\epsilon^{(n)}(t),
\end{equation}
where \(I_{M^{(n)}}\) is the \(M^{(n)}\times M^{(n)}\) identity matrix and
\begin{equation}
\boldsymbol{\chi}_\epsilon^{(n)}(t)= \begin{pmatrix}
\chi_{\epsilon_1}^{(n)}(t) \\
\chi_{\epsilon_2}^{(n)}(t) \\
\vdots \\
\chi_{\epsilon_{M^{(n)}}}^{(n)}(t)
\end{pmatrix}.
\end{equation}

Stacking all pulsars, the full continuous model is
\begin{equation}
\frac{d}{dt}X = \mathcal{A}\,X + \Xi(t),
\end{equation}
with
\begin{equation}
\mathcal{A} = \operatorname{diag}\Bigl(0_{M^{(1)}\times M^{(1)}},\,0_{M^{(2)}\times M^{(2)}},\,\dots,\,0_{M^{(N)}\times M^{(N)}}\Bigr),
\end{equation}
and
\begin{equation}
\Xi(t)= \begin{pmatrix}
I_{M^{(1)}}\,\boldsymbol{\chi}_\epsilon^{(1)}(t) \\
I_{M^{(2)}}\,\boldsymbol{\chi}_\epsilon^{(2)}(t) \\
\vdots \\
I_{M^{(N)}}\,\boldsymbol{\chi}_\epsilon^{(N)}(t)
\end{pmatrix}.
\end{equation}

Since \(\mathcal{A}=0\), the solution is given solely by the noise integration. For a discrete time step \(\Delta t\) the exact discretisation for pulsar \(n\) is
\begin{equation}
\delta\boldsymbol{\epsilon}^{(n)}_{k+1} = e^{0\cdot \Delta t}\,\delta\boldsymbol{\epsilon}^{(n)}_k + \eta_k^{(n)} = I_{M^{(n)}}\,\delta\boldsymbol{\epsilon}^{(n)}_k + \eta_k^{(n)},
\end{equation}
with the noise increment defined as
\begin{equation}
\eta_k^{(n)} = \int_0^{\Delta t}\boldsymbol{\chi}_\epsilon^{(n)}(t_k+s)\,ds.
\end{equation}

Thus, the state--transition matrix for pulsar \(n\) is
\begin{equation}
F^{(n)} = I_{M^{(n)}}.
\end{equation}

Assuming that the noise processes are white with
\begin{equation}
\langle \chi_{\epsilon_m}^{(n)}(t)\,\chi_{\epsilon_{m'}}^{(n')}(t') \rangle 
= \left[\sigma_{\epsilon_m}^{(n)}\right]^2\,\delta_{n,n'}\,\delta_{m,m'}\,\delta(t-t'),
\end{equation}
the covariance of the noise increment for pulsar \(n\) is
\begin{equation}
Q^{(n)} = \langle \eta_k^{(n)}\,\eta_k^{(n)\top} \rangle 
= \Delta t\,\operatorname{diag}\Bigl(\left[\sigma_{\epsilon_1}^{(n)}\right]^2,\left[\sigma_{\epsilon_2}^{(n)}\right]^2,\dots,\left[\sigma_{\epsilon_{M^{(n)}}}^{(n)}\right]^2\Bigr).
\end{equation}

Stacking all pulsars, the full discrete model becomes
\begin{equation}
X_{k+1} = F\,X_k + \eta_k,
\end{equation}
with the overall state--transition matrix
\begin{equation}
F = \operatorname{diag}\Bigl(F^{(1)},\,F^{(2)},\,\dots,\,F^{(N)}\Bigr)
= \operatorname{diag}\Bigl(I_{M^{(1)}},\,I_{M^{(2)}},\,\dots,\,I_{M^{(N)}}\Bigr),
\end{equation}
and the full noise covariance
\begin{equation}
Q = \operatorname{diag}\Bigl(Q^{(1)},\,Q^{(2)},\,\dots,\,Q^{(N)}\Bigr).
\end{equation}

\bigskip
\textbf{Summary:}
\begin{itemize}
    \item For each pulsar \(n\) with \(M^{(n)}\) parameters, the discrete state--transition matrix is
    \begin{equation}
    F^{(n)} = I_{M^{(n)}}.
    \end{equation}
    \item The process noise covariance is
    \begin{equation}
    Q^{(n)} = \Delta t\,\operatorname{diag}\Bigl(\left[\sigma_{\epsilon_1}^{(n)}\right]^2,\left[\sigma_{\epsilon_2}^{(n)}\right]^2,\dots,\left[\sigma_{\epsilon_{M^{(n)}}}^{(n)}\right]^2\Bigr).
    \end{equation}
    \item The overall discrete model is
    \begin{equation}
    X_{k+1} = F\,X_k + \eta_k, \quad \text{with} \quad F = \operatorname{diag}\Bigl(I_{M^{(1)}},\,I_{M^{(2)}},\,\dots,\,I_{M^{(N)}}\Bigr),
    \end{equation}
    and
    \begin{equation}
    Q = \operatorname{diag}\Bigl(Q^{(1)},\,Q^{(2)},\,\dots,\,Q^{(N)}\Bigr).
    \end{equation}
\end{itemize}




\section{Kalman filter implementation}



\section{Predict step}
Assume the state is ordered as
\begin{equation}
x = \begin{pmatrix} x_{GW} \\[1mm] x_{\phi/f} \\[1mm] x_{\delta\epsilon} \end{pmatrix},
\end{equation}
and that after the measurement update the covariance is
\begin{equation}
P^+ = \begin{pmatrix}
P^+_{GW} & P^+_{GW,\phi/f} & P^+_{GW,\delta\epsilon}\\[1mm]
P^+_{\phi/f,GW} & P^+_{\phi/f} & P^+_{\phi/f,\delta\epsilon}\\[1mm]
P^+_{\delta\epsilon,GW} & P^+_{\delta\epsilon,\phi/f} & P^+_{\delta\epsilon}
\end{pmatrix}.
\end{equation}
Suppose also that the state transition matrix and process noise are block diagonal:
\begin{equation}
F = \begin{pmatrix}
F_{GW} & 0 & 0\\[1mm]
0 & F_{\phi/f} & 0\\[1mm]
0 & 0 & F_{\delta\epsilon}
\end{pmatrix},\qquad
Q = \begin{pmatrix}
Q_{GW} & 0 & 0\\[1mm]
0 & Q_{\phi/f} & 0\\[1mm]
0 & 0 & Q_{\delta\epsilon}
\end{pmatrix}.
\end{equation}
Then the prediction step is given by
\begin{equation}
P_p = F\,P^+\,F^T + Q,
\end{equation}
which expands to
\begin{equation}
P_p = \begin{pmatrix}
F_{GW}\,P^+_{GW}\,F_{GW}^T + Q_{GW} & F_{GW}\,P^+_{GW,\phi/f}\,F_{\phi/f}^T & F_{GW}\,P^+_{GW,\delta\epsilon}\,F_{\delta\epsilon}^T\\[1mm]
F_{\phi/f}\,P^+_{\phi/f,GW}\,F_{GW}^T & F_{\phi/f}\,P^+_{\phi/f}\,F_{\phi/f}^T + Q_{\phi/f} & F_{\phi/f}\,P^+_{\phi/f,\delta\epsilon}\,F_{\delta\epsilon}^T\\[1mm]
F_{\delta\epsilon}\,P^+_{\delta\epsilon,GW}\,F_{GW}^T & F_{\delta\epsilon}\,P^+_{\delta\epsilon,\phi/f}\,F_{\phi/f}^T & F_{\delta\epsilon}\,P^+_{\delta\epsilon}\,F_{\delta\epsilon}^T + Q_{\delta\epsilon}
\end{pmatrix}.
\end{equation}
Thus:
\begin{itemize}
    \item The GW block is 
    \begin{equation}
    P_{p}^{GW} = F_{GW}\,P^+_{GW}\,F_{GW}^T + Q_{GW}.
    \end{equation}
    \item The \(\phi/f\) block is 
    \begin{equation}
    P_{p}^{\phi/f} = F_{\phi/f}\,P^+_{\phi/f}\,F_{\phi/f}^T + Q_{\phi/f}.
    \end{equation}
    \item The \(\delta\epsilon\) block is 
    \begin{equation}
    P_{p}^{\delta\epsilon} = F_{\delta\epsilon}\,P^+_{\delta\epsilon}\,F_{\delta\epsilon}^T + Q_{\delta\epsilon}.
    \end{equation}
    \item The cross–covariance between GW and \(\phi/f\) is 
    \begin{equation}
    P_{p}^{GW,\phi/f} = F_{GW}\,P^+_{GW,\phi/f}\,F_{\phi/f}^T.
    \end{equation}
    \item The cross–covariance between GW and \(\delta\epsilon\) is 
    \begin{equation}
    P_{p}^{GW,\delta\epsilon} = F_{GW}\,P^+_{GW,\delta\epsilon}\,F_{\delta\epsilon}^T.
    \end{equation}
    \item The cross–covariance between \(\phi/f\) and \(\delta\epsilon\) is 
    \begin{equation}
    P_{p}^{\phi/f,\delta\epsilon} = F_{\phi/f}\,P^+_{\phi/f,\delta\epsilon}\,F_{\delta\epsilon}^T.
    \end{equation}
\end{itemize}




Assume that the state is partitioned into three sectors as
\begin{equation}
x=\begin{pmatrix}
x_{GW}\\[1mm]
x_{\phi/f}\\[1mm]
x_{\delta\epsilon}
\end{pmatrix},
\end{equation}
with prior (predicted) covariance
\begin{equation}
P^-=\begin{pmatrix}
P^-_{GW} & 0 & 0\\[1mm]
0 & P^-_{\phi/f} & 0\\[1mm]
0 & 0 & P^-_{\delta\epsilon}
\end{pmatrix}.
\end{equation}
Suppose that the measurement equation for the currently observed pulsar (say pulsar \(n\)) is
\begin{equation}
y=\frac{1}{f_0^{(n)}}\,\delta\phi^{(n)}+\mathbf{M}^{(n)}\,\boldsymbol{\delta\epsilon}^{(n)}-r^{(n)}.
\end{equation}
Then the measurement matrix is written in block form as
\begin{equation}
H=\begin{pmatrix} H_{GW} & H_{\phi/f} & H_{\delta\epsilon} \end{pmatrix},
\end{equation}
where
\begin{equation}
H_{GW}=[0\,\cdots\,0\,\, -1\, 0\,\cdots\,0] \quad,\quad
H_{\phi/f}=[0\,\cdots\,0\,\, 1/f_0^{(n)}\, 0\,\cdots\,0],
\end{equation}
\begin{equation}
H_{\delta\epsilon}=[0\,\cdots\,0\,\, \mathbf{M}^{(n)}\, 0\,\cdots\,0].
\end{equation}
Define
\begin{equation}
S=H_{GW}P^-_{GW}H_{GW}^T+H_{\phi/f}P^-_{\phi/f}H_{\phi/f}^T+H_{\delta\epsilon}P^-_{\delta\epsilon}H_{\delta\epsilon}^T+R.
\end{equation}
Then the Kalman gain is
\begin{equation}
K=P^-H^T\,S^{-1},
\end{equation}
which in block form is
\begin{equation}
K=\begin{pmatrix}
K_{GW}\\[1mm]
K_{\phi/f}\\[1mm]
K_{\delta\epsilon}
\end{pmatrix}=
\begin{pmatrix}
P^-_{GW}H_{GW}^T\\[1mm]
P^-_{\phi/f}H_{\phi/f}^T\\[1mm]
P^-_{\delta\epsilon}H_{\delta\epsilon}^T
\end{pmatrix}\,S^{-1}.
\end{equation}

The state update is
\begin{equation}
x^+=x^-+K\Bigl(y-Hx^-\Bigr).
\end{equation}

The covariance update is
\begin{equation}
P^+=(I-KH)P^-.
\end{equation}
Since \(P^-\) is block diagonal, this expands to
\begin{equation}
P^+ = \begin{pmatrix}
P^-_{GW}-K_{GW}H_{GW}P^-_{GW} & -K_{GW}H_{\phi/f}P^-_{\phi/f} & -K_{GW}H_{\delta\epsilon}P^-_{\delta\epsilon}\\[1mm]
-K_{\phi/f}H_{GW}P^-_{GW} & P^-_{\phi/f}-K_{\phi/f}H_{\phi/f}P^-_{\phi/f} & -K_{\phi/f}H_{\delta\epsilon}P^-_{\delta\epsilon}\\[1mm]
-K_{\delta\epsilon}H_{GW}P^-_{GW} & -K_{\delta\epsilon}H_{\phi/f}P^-_{\phi/f} & P^-_{\delta\epsilon}-K_{\delta\epsilon}H_{\delta\epsilon}P^-_{\delta\epsilon}
\end{pmatrix}.
\end{equation}

Assume that the state is partitioned into three sectors:
\begin{equation}
x=\begin{pmatrix}
x_{GW}\\[1mm]
x_{\phi/f}\\[1mm]
x_{\delta\epsilon}
\end{pmatrix},
\end{equation}
and that the measurement equation for pulsar \(n\) is
\begin{equation}
y=\frac{1}{f_0^{(n)}}\,\delta\phi^{(n)}+\mathbf{M}^{(n)}\,\boldsymbol{\delta\epsilon}^{(n)}-r^{(n)}.
\end{equation}
In our state ordering, the measurement matrix has nonzero blocks only in the following positions:
\begin{itemize}
    \item In the GW sector, a \(-1\) in the entry corresponding to \(r^{(n)}\);
    \item In the \(\phi/f\) sector, a \(1/f_0^{(n)}\) in the entry corresponding to \(\delta\phi^{(n)}\);
    \item In the \(\delta\epsilon\) sector, the row vector \(\mathbf{M}^{(n)}\) in the block corresponding to pulsar \(n\).
\end{itemize}
Let
\begin{equation}
S=P^-_{r^{(n)}r^{(n)}}+\frac{1}{f_0^{(n)2}}P^-_{\delta\phi^{(n)}\delta\phi^{(n)}}+\mathbf{M}^{(n)}P^-_{\delta\epsilon^{(n)}\delta\epsilon^{(n)}}\mathbf{M}^{(n)T}+R.
\end{equation}
Then the Kalman gains for the measured components are
\begin{align}
K_{r^{(n)}} &= -\frac{P^-_{r^{(n)}r^{(n)}}}{S},\\[1mm]
K_{\delta\phi^{(n)}} &= \frac{P^-_{\delta\phi^{(n)}\delta\phi^{(n)}}}{f_0^{(n)}\,S},\\[1mm]
K_{\delta\epsilon^{(n)}} &= \frac{P^-_{\delta\epsilon^{(n)}\delta\epsilon^{(n)}}\,\mathbf{M}^{(n)T}}{S}.
\end{align}
All other entries of the gain \(K\) are zero.

The covariance update is
\begin{equation}
P^+=P^--K\,H\,P^-.
\end{equation}
Thus, the updated variances for the measured components are
\begin{align}
P^+_{r^{(n)}r^{(n)}} &= P^-_{r^{(n)}r^{(n)}}-\frac{\Bigl(P^-_{r^{(n)}r^{(n)}}\Bigr)^2}{S},\\[1mm]
P^+_{\delta\phi^{(n)}\delta\phi^{(n)}} &= P^-_{\delta\phi^{(n)}\delta\phi^{(n)}}-\frac{\Bigl(P^-_{\delta\phi^{(n)}\delta\phi^{(n)}/f_0^{(n)}}\Bigr)^2}{S},\\[1mm]
P^+_{\delta\epsilon^{(n)}\delta\epsilon^{(n)}} &= P^-_{\delta\epsilon^{(n)}\delta\epsilon^{(n)}}-\frac{P^-_{\delta\epsilon^{(n)}\delta\epsilon^{(n)}}\,\mathbf{M}^{(n)T}\,\mathbf{M}^{(n)}\,P^-_{\delta\epsilon^{(n)}\delta\epsilon^{(n)}}}{S}.
\end{align}
Similarly, the cross–covariances are updated as
\begin{align}
P^+_{r^{(n)},\delta\phi^{(n)}} &= P^-_{r^{(n)},\delta\phi^{(n)}}+\frac{P^-_{r^{(n)}r^{(n)}}\,P^-_{\delta\phi^{(n)}\delta\phi^{(n)}/f_0^{(n)}}}{S},\\[1mm]
P^+_{r^{(n)},\delta\epsilon^{(n)}} &= P^-_{r^{(n)},\delta\epsilon^{(n)}}+\frac{P^-_{r^{(n)}r^{(n)}}\,P^-_{\delta\epsilon^{(n)}\delta\epsilon^{(n)}}\,\mathbf{M}^{(n)T}}{S},\\[1mm]
P^+_{\delta\phi^{(n)},\delta\epsilon^{(n)}} &= P^-_{\delta\phi^{(n)},\delta\epsilon^{(n)}}-\frac{P^-_{\delta\phi^{(n)}\delta\phi^{(n)}/f_0^{(n)}}\,P^-_{\delta\epsilon^{(n)}\delta\epsilon^{(n)}}\,\mathbf{M}^{(n)T}}{S}.
\end{align}
All other blocks of \(P^+\) remain unchanged.

Finally, the prediction step is
\begin{equation}
P_p = F\,P^+\,F^T+Q,
\end{equation}
and since \(F\) and \(Q\) are block diagonal,
\begin{equation}
P_p = \begin{pmatrix}
F_{GW}\,P^+_{GW}\,F_{GW}^T+Q_{GW} & F_{GW}\,P^+_{GW,\phi/f}\,F_{\phi/f}^T & F_{GW}\,P^+_{GW,\delta\epsilon}\,F_{\delta\epsilon}^T\\[1mm]
F_{\phi/f}\,P^+_{\phi/f,GW}\,F_{GW}^T & F_{\phi/f}\,P^+_{\phi/f}\,F_{\phi/f}^T+Q_{\phi/f} & F_{\phi/f}\,P^+_{\phi/f,\delta\epsilon}\,F_{\delta\epsilon}^T\\[1mm]
F_{\delta\epsilon}\,P^+_{\delta\epsilon,GW}\,F_{GW}^T & F_{\delta\epsilon}\,P^+_{\delta\epsilon,\phi/f}\,F_{\phi/f}^T & F_{\delta\epsilon}\,P^+_{\delta\epsilon}\,F_{\delta\epsilon}^T+Q_{\delta\epsilon}
\end{pmatrix}.
\end{equation}






\newpage
\appendix 

**Note: everything below here is a scratch space.** 
\section{Derivation of the Measurement Equation}

\subsection{Introduction}
We seek to apply the Kalman filter method to real PTA data, rather than a measured frequency time series $f_{\textrm{m}}^{(n)}(t)$, as done in previous work. \newline 

\noindent By ``real data," we refer to a \textsc{.tim} file (which contains the TOAs) and a \textsc{.par} file (which provides constrained \textit{a priori} estimates of the pulsar parameters, such as sky position, spin frequency, etc.). \newline

For the purposes of this note, we assume that the \textsc{.tim} and \textsc{.par} files have been processed through a standard pulsar timing library (e.g., TEMPO or PINT) to produce timing residuals $\delta t$. These timing residuals define our measurement vector, $\boldsymbol{Y}$. \newline 

Pulsar timing libraries employ a deterministic, parameterized model that utilizes estimates of the timing ephemeris parameters, $\boldsymbol{\hat{\theta}}$ , to predict the pulse TOA, $t_{\text{det}}(\boldsymbol{\hat{\theta}})$. This deterministic model accounts for standard effects such as Shapiro delay, proper motions, and binary interactions.\footnote{One might reasonably question whether removing all known contributions first, before analyzing residuals, could inadvertently filter out part of the signal of interest. This potential issue is well known in the pulsar community and will be addressed later.} The timing residuals are then defined as the difference between the actual TOAs and the predicted TOAs:

\begin{equation}
	\delta t = t_{\text{TOA}} - t_{\text{det}}(\boldsymbol{\hat{\theta}}). \label{eq:delta_t_og_defn}
\end{equation}


\subsection{Definition of phases, $\phi(t)$}

In previous works, we have a hidden state $\boldsymbol{X}$ which we identify with the intrinsic pulsar frequency $f_p(t)$ and a measurement  $\boldsymbol{Y}$ which we identify with a measured frequency $f_m(t)$. The two are related via a measurement equation 
\begin{equation}
	f_{\text{m}}(t) = f_{\text{p}}(t) \left[1 - a(t)\right], \label{eq:measurement}
\end{equation}
where $a(t)$ quantifies the influence of the GW (for simplicity, we omit here the superscript-$(n)$ notation.) \newline 

\noindent We want to derive an equivalent measurement equation that relates the new measurement $\boldsymbol{Y} = \delta t$ with some hidden states $\boldsymbol{X}$. \newline 


\noindent To start, separate the intrinsic pulsar frequency into deterministic and stochastic parts $f_{\text{p}}(t) = \bar{f}(t) + \delta f(t)$ and define the following phase variables:
\begin{align}
	\text{Measured phase:} \quad & \phi_{\text{m}}(t) = \int_0^t f_{\text{m}}(t') dt', \label{eq:measured_phase}\\
	\text{Intrinsic phase:} \quad & \phi_{\text{p}}(t;\boldsymbol{\theta}) = \int_0^t \bar{f}(t'; \boldsymbol{\theta}) dt', \label{eq:model_phase} \\
	\text{Model phase:} \quad & \phi_{\text{p}}(t;\boldsymbol{\hat{\theta}}) = \int_0^t \bar{f}(t'; \boldsymbol{\hat{\theta}}) dt'.
\end{align}

Here, $\boldsymbol{\theta}$ represents the true timing-ephemeris parameters, while $\boldsymbol{\hat{\theta}}$ are the best-fit model estimates (i.e., the true parameters satisfy $\boldsymbol{\theta} = \boldsymbol{\hat{\theta}} + \delta \boldsymbol{\theta}$). In a standard PTA analysis it is assumed that any difference between the true determinisic solution and the estimated solution will be small (see e.g. \href{https://arxiv.org/abs/2105.13270}{Taylor 2021, Section 7.1}) and we can therefore linearize as
\begin{eqnarray}
\phi_{\text{p}}(t;\boldsymbol{\theta}) = \phi_{\text{p}}(t;\boldsymbol{\hat{\theta}}) + \mathbf{M}_{\phi} \delta \boldsymbol{\theta} \label{eq:MMatrix}
\end{eqnarray}
where $\mathbf{M}_{\phi}$ is the design matrix of partial derivatives (of the phases) with respect to the parameters i.e. $= \partial_{\theta} \phi(t)$. 

\subsection{Definition of the timing residual, $\delta t$}
The timing residual is given by Equation \eqref{eq:delta_t_og_defn}. We can also provide an equivalent definition in terms of phases / frequencies as follows. \newline 

\noindent According to the timing-ephemeris model, the $n$-th pulse arrives at time $t$ if 

\begin{equation}
	\phi_{\text{p}}(t; \boldsymbol{\hat{\theta}}) = n.
	\end{equation}
for some integer $n$. The actual pulse arrives at time $t + \delta t$, satisfying:

\begin{equation}
	\phi_{\text{m}}(t + \delta t) = n.
\end{equation}

\noindent A linear expansion gives
\begin{align}
	\phi_{\text{m}}(t + \delta t)  &\approx \phi_{\text{m}}(t) + \frac{d}{dt} \phi_{\text{m}}(t) \delta t \nonumber \\
	&= \phi_{\text{m}}(t) + f_{\text{m}}(t) \delta t. \label{eq:assumption 1}
\end{align}
Rearranging, we express the timing residual in terms of phases:
\begin{equation}
	\delta t (t)= \frac{\phi_{\text{p}}(t; \boldsymbol{\hat{\theta}}) - \phi_{\text{m}}(t)}{f_{\text{m}}(t)}. \label{eq:delta_t}
\end{equation}


\subsection{Expanding $\phi_m(t)$}
From Equation \eqref{eq:measurement} and Equation \eqref{eq:measured_phase}, we can express the measured phase as
\begin{align}
	\phi_m(t)
	&= \int_{0}^{t} f_m(\tau)\,d\tau \\
	&= \int_{0}^{t} 
	\bigl[1 - a(\tau)\bigr]\,
	\bigl[\bar{f}(\tau) + \delta f(\tau)\bigr]\,d\tau \\
	&= \underbrace{\int_{0}^{t} \bar{f}(\tau)\,d\tau}_{\text{timing solution}}
	\;+\;
	\underbrace{\int_{0}^{t} \delta f(\tau)\,d\tau}_{\text{red noise}}
	\;-\;
	\underbrace{\int_{0}^{t} a(\tau)\,\bar{f}(\tau)\,d\tau}_{\text{GW term}}
	\;-\;
	\underbrace{\int_{0}^{t} a(\tau)\,\delta f(\tau)\,d\tau}_{\text{small term}}. \label{eq:phi_m_decomposed}
\end{align}



We make the following observations
\begin{itemize}
	\item The timing solution term is just the phase 	$\phi_{\text{p}}(t;\boldsymbol{\theta})$, Equation \eqref{eq:model_phase}
	\item The red noise term integral is just a phase, a fluctuation from the determinisic timing model, which we will call $\delta \phi$.
	\item Because the spin-down derivative term is small, the determinstic part of the GW term can be approximated as constant, $\bar{f}(t) = f_0 + \dot{f} t \approx f_0$. We will define $r(t) =\int_{0}^{t} a(\tau) \,d\tau $, (see e.g. \href{https://arxiv.org/abs/1003.0677}{Equation 5 of Sesana \& Vecchio 2010})
	\item The second order term is small and can be dropped.
\end{itemize}
We can therefore write Equation \eqref{eq:phi_m_decomposed} more concisely as 
\begin{equation}
	\phi_m(t)
	= \phi_{\text{p}}(t;\boldsymbol{\theta})
	\;+\;
	\delta \phi(t)
	\;-\;
	f_0 r(t) \label{eq:phi_m_concise}
\end{equation}

\subsection{Putting it all together}
Combining Equations \eqref{eq:MMatrix}, \eqref{eq:delta_t}, and \eqref{eq:phi_m_concise}, and taking the denominator of Equation \eqref{eq:delta_t} to be constant ($=f_0$) we can write
\begin{align}
	\delta t &= \frac{1}{f_0} \left[ \phi_{\text{p}}(t;\boldsymbol{\hat{\theta}}) - \phi_{\text{p}}(t;\boldsymbol{\theta}) - \delta \phi(t) + f_0 r(t)\right] \\
	&= - \frac{1}{f_0}M_{\phi} \delta \boldsymbol{\theta} - \frac{\delta \phi}{f_0} + r(t)
\end{align}
In practice, TEMPO/PINT return a design matrix defined in terms of partial derivatives of the TOAs, rather than in terms of phases, so we can just write
\begin{equation}
	\boxed{\delta t = \boldsymbol{M}_{\text{TOA}} \delta \boldsymbol{\theta} - \frac{\delta \phi}{f_0} + r(t)} \label{eq:measurement_new}
\end{equation}
Equation \eqref{eq:measurement_new} is the measurement equation. \newline 

\textbf{Summary of assumptions and approximations }
\begin{itemize}
	\item  Linearisation to define the M-matrix in Equation \eqref{eq:MMatrix}
	\item Linear expansion and neglecting higher-order terms in Equation \eqref{eq:assumption 1}.
	\item Approximating $\bar{f}(t) \approx f_0$ in the GW term of Equation \eqref{eq:phi_m_decomposed}.
	\item Approximating $f_{\text{m}} \approx f_0$ in the denominator of Equation \eqref{eq:delta_t}.
	\item Neglecting small second-order terms in Equation \eqref{eq:phi_m_decomposed}.

\end{itemize}





\section{Single pulsar}\label{sec:single_pulsar}
To start, lets put the above into a state-space frame work with a single pulsar, $N=1$. \newline 

\subsection{State vector}

\noindent The state vector is 

\begin{equation}
	\boldsymbol{X} = \left(\delta \phi, \delta f, r,a,\delta \epsilon_1, \delta \epsilon_2, \cdots, \delta \epsilon_{\mathrm M} \right)
\end{equation}
The first two variables, $\delta \phi$ and $\delta f$, are deviations from the spin-down parameters in timing model fit. Since we know these deviations won't move us more than 1 turn away (because these are nice millisecond pulsars) these are reasonable state variables to use. So for example, $\delta f$ is a fluctuation from the measured spin-down in the timing model. \newline 


\noindent The subsequent two variables, $r$ and $a$, quantify the effect of the gravitational wave. The variable $a(t)$ is the familiar redshift quantity from \href{https://arxiv.org/abs/2501.06990}{PTA P3}, and $r$ is the induced timing residual, obtained by integrating the redshift (see e.g. \href{https://arxiv.org/abs/1003.0677}{Equation 5 of Sesana \& vecchio 2010}). \newline 


\noindent The final $M$ parameters, $\delta \epsilon_i$ are parameters of the design matrix. These parameters will be used to incorporate the effects of the timing model. \newline 



\noindent The observation vector is 
\begin{equation}
	\boldsymbol{Y} = \left(\delta t \right)
\end{equation}
where $\delta t$ is the timing residual.


\subsection{Dynamical Equations (continuous time)}

The above state variables evolve according to the following dynamical equations\footnote{TK: these might not be the equations to use. For instance, do we still need a $\gamma$ term? I just use these as place holders for now.}


\begin{align}
	\frac{d}{dt} \delta \phi^{(n)} &= \delta f^{(n)}, \\
	\frac{d}{dt} \delta f^{(n)} &= -\gamma_p^{(n)} \delta f^{(n)} + \chi_p^{(n)}(t;\sigma_p^{(n)}), \\
	\frac{d}{dt} r^{(n)} &= a^{(n)}, \\
	\frac{d}{dt} a^{(n)} &= -\gamma_a a^{(n)} + \chi_a^{(n)}(t;\sigma_a^{(n)}), \\
	\frac{d}{dt} \delta \epsilon^{(n)}_{m} &= \chi_\epsilon^{(n)}(t;\sigma_\epsilon) \quad \forall m \in [1, M^{(n)}].
\end{align}
where $\chi_a$ is the usual white noise stochastic process. \newline 



\noindent The states are related to the observables via a measurement equation
\begin{equation}
	\delta t = \frac{\delta \phi}{f_0} + \mathbf{M} \mathbf{\delta \epsilon} - r
\end{equation}
where  $\mathbf{M}$ is the \textit{design matrix} and $f_0$ is the pulsar rotation frequency \footnote{I need to check the sign of the residual term. Is it definitely a minus?}. \newline 


\noindent Note that the dynamics of the $\delta \epsilon^{(n)}_{m}$ terms here are a bit different to how it is done in \textsc{minnow}. In \textsc{minnow} these terms have zero process noise, and they just get initialised with some values (and covariance) which can then be estimated by the inference procedure. In our case I think we can marginalise over these timing model parameters by including them in the state and setting the variance $\sigma_\epsilon$ to be very large. This is effecively what pulsar people already do when using Gaussian processes in \textsc{enterprise}; see e.g. \href{https://arxiv.org/abs/1407.1838}{Section IV of van Haasteran \& Vallisneri, 2014}. In this way, we can avoid having to estimate an extra $\sum_{i=1}^{N} M^{(i)}$ parameters. I include the process noise term here by default; if we want to remove it in the future and do the full parameter estimation, we can just set $\sigma_\epsilon=0$ everywhere that follows.




\subsection{Discretisation}

\subsubsection{F-matrix}
	
	For the $\left(\delta\phi,\delta f\right)$ block, the continuous system is
	\begin{equation}
	\frac{d}{dt}\begin{pmatrix}\delta\phi \\ \delta f\end{pmatrix} 
	=
	\begin{pmatrix}
		0 & 1\\[1mm]
		0 & -\gamma_p
	\end{pmatrix}
	\begin{pmatrix}\delta\phi \\ \delta f\end{pmatrix}
	+
	\begin{pmatrix}0\\1\end{pmatrix}\chi_p(t),
	\end{equation}
	and the exact discretisation (matrix exponential) gives
	\begin{equation}
	F_p = \begin{pmatrix}
		1 & \dfrac{1-e^{-\gamma_p\Delta t}}{\gamma_p}\\[1mm]
		0 & e^{-\gamma_p\Delta t}
	\end{pmatrix}.
	\end{equation}
	Similarly, for the $\left(r,a\right)$ block we have
	\begin{equation}
	F_a = \begin{pmatrix}
		1 & \dfrac{1-e^{-\gamma_a\Delta t}}{\gamma_a}\\[1mm]
		0 & e^{-\gamma_a\Delta t}
	\end{pmatrix}.
	\end{equation}
For the timing model parameters $\delta\epsilon_m$, the dynamics are a random walk:
	\begin{equation}
	\delta\epsilon_{m,k+1} = \delta\epsilon_{m,k} + \eta_{\epsilon,m,k},
	\end{equation}
	so that the deterministic evolution is simply the identity. \newline 
	
	\noindent Thus, the overall discrete state transition matrix is given by the block diagonal matrix
	\begin{equation}
	\boldsymbol{F} = \begin{pmatrix}
		F_p & 0 & 0\\[1mm]
		0 & F_a & 0\\[1mm]
		0 & 0 & I_M
	\end{pmatrix},
	\end{equation}

	
\subsubsection{Q-matrix}
	
	For the $\left(\delta\phi,\delta f\right)$ block with noise variance $\sigma_p^2$, the discretised noise covariance is
	\begin{equation}
	Q_p = \sigma_p^2
	\begin{pmatrix}
		\displaystyle \frac{\Delta t}{\gamma_p^2} - \frac{2\left(1-e^{-\gamma_p\Delta t}\right)}{\gamma_p^3} + \frac{1-e^{-2\gamma_p\Delta t}}{2\gamma_p^3} & \displaystyle \frac{1-e^{-\gamma_p\Delta t}}{\gamma_p^2} - \frac{1-e^{-2\gamma_p\Delta t}}{2\gamma_p^2}\\[2mm]
		\displaystyle \frac{1-e^{-\gamma_p\Delta t}}{\gamma_p^2} - \frac{1-e^{-2\gamma_p\Delta t}}{2\gamma_p^2} & \displaystyle \frac{1-e^{-2\gamma_p\Delta t}}{2\gamma_p}
	\end{pmatrix}.
	\end{equation}
	For the $\left(r,a\right)$ block with noise variance $\sigma_a^2$, we similarly have
	\begin{equation}
	Q_a = \sigma_a^2
	\begin{pmatrix}
		\displaystyle \frac{\Delta t}{\gamma_a^2} - \frac{2\left(1-e^{-\gamma_a\Delta t}\right)}{\gamma_a^3} + \frac{1-e^{-2\gamma_a\Delta t}}{2\gamma_a^3} & \displaystyle \frac{1-e^{-\gamma_a\Delta t}}{\gamma_a^2} - \frac{1-e^{-2\gamma_a\Delta t}}{2\gamma_a^2}\\[2mm]
		\displaystyle \frac{1-e^{-\gamma_a\Delta t}}{\gamma_a^2} - \frac{1-e^{-2\gamma_a\Delta t}}{2\gamma_a^2} & \displaystyle \frac{1-e^{-2\gamma_a\Delta t}}{2\gamma_a}
	\end{pmatrix}.
	\end{equation}
	For each timing model parameter $\delta\epsilon_m$, modeled as a random walk,
	\begin{equation}
	\delta\epsilon_{m,k+1} = \delta\epsilon_{m,k} + \eta_{\epsilon,m,k}, \quad \eta_{\epsilon,m,k}\sim \mathcal{N}(0,\sigma_\epsilon^2\,\Delta t),
	\end{equation}
	so that
	\begin{equation}
	Q_\epsilon = \sigma_\epsilon^2\,\Delta t\,I_M.
	\end{equation}
	Thus, the overall process noise covariance is block diagonal:
	\begin{equation}
	\boldsymbol{Q} = \begin{pmatrix}
		Q_p & 0 & 0\\[2mm]
		0 & Q_a & 0\\[2mm]
		0 & 0 & \sigma_\epsilon^2\,\Delta t\,I_M
	\end{pmatrix}.
	\end{equation}
	
\subsubsection{H-matrix}
	Recall that the measurement equation is
		\begin{equation}
	\delta t = \frac{\delta\phi}{f_0} + \boldsymbol{M}\,\boldsymbol{\delta\epsilon} - r.
		\end{equation}
	Since the state vector is
	\begin{equation}
	\boldsymbol{X} = \begin{pmatrix}
		\delta\phi \\ \delta f \\ r \\ a \\ \delta\epsilon_1 \\ \vdots \\ \delta\epsilon_M
	\end{pmatrix},
		\end{equation}
	the measurement depends only on $\delta\phi$, $r$, and the $\delta\epsilon_m$. Therefore, the measurement matrix is
		\begin{equation}
	\boldsymbol{H} = \begin{pmatrix}
		\frac{1}{f_0} & 0 & -1 & 0 & M_1 & \cdots & M_M
	\end{pmatrix},
		\end{equation}
	so that
		\begin{equation}
	\delta t = \boldsymbol{H}\,\boldsymbol{X}.
		\end{equation}
	
\subsubsection{R-matrix}
	Assuming the measurement noise is white with variance $\sigma_t^2$, and given that there is a single measurement per time step, we have
	\begin{equation}
	\boldsymbol{R} = \sigma_t^2.
	\end{equation}

	


\section{Multiple pulsars}

The formulation in Section \ref{sec:single_pulsar} extends straightforwardly to multiple pulsars. There are two complications which must be handled with care.

\begin{enumerate}
	\item The covariance between the $a^{(n)}$ terms for different pulsars. This is the Hellings-Downs effect.
	\item The fact that in general all pulsars are observed at different times. This means that instead of having a nice observation vector $\mathbf{Y}$ of length $N$, in general we just have a single observation from a single pulsar at a given timestep. 
\end{enumerate}


\noindent Regarding (1), the ensemble statistics of $\chi_{\mathrm a}^{(n)}$ are
\begin{align}
	\langle \chi^{(n)}_{\mathrm a}(t) \rangle &= 0 \ , 	\label{eq:xieqn1} \\
	\langle \chi^{(n)}_{\mathrm a}(t) \, \chi^{(n')}_{\mathrm a}(t') \rangle &= \left[\sigma^{(n,n')}_{\mathrm a}\right]^2 \delta(t - t') \ .	\label{eq:xieqn2}
\end{align}
with
\begin{eqnarray}
	\left[\sigma^{(n,n')}_{\mathrm a}\right]^2 = \frac{\langle h^2\rangle}{6} \gamma_{\mathrm a} \, \Gamma \left[ \theta^{(n,n')} \right] \, , \label{eq:sigma_a_expression}
\end{eqnarray}
where $\langle h^2 \rangle$ is the mean square GW strain from the $M$ summed sources at the Earth's position, $\theta^{(n,n')}$ is the angle between the $n$-th and $n'$-th pulsars, and one has
\begin{eqnarray}
	\Gamma\left[\theta^{(n,n')} \right] =  \frac{3}{2} x_{n n'} \ln x_{n n'}  -\frac{x_{n n'} }{4}+\frac{1}{2} + \frac{1}{2} \delta_{n n'}\label{eq:correlation} \, ,
\end{eqnarray}
with $x_{nn'} = \left[1 - \cos \theta^{(n,n')}\right]/2$. 



\subsection{Discretisation for multiple pulsars}
	
	For $N$ pulsars, we define the stacked state vector
	\begin{equation}
		\boldsymbol{X} = \begin{pmatrix}
			\boldsymbol{X}^{(1)} \\
			\boldsymbol{X}^{(2)} \\
			\vdots \\
			\boldsymbol{X}^{(N)}
		\end{pmatrix}, \qquad \text{with} \qquad
		\boldsymbol{X}^{(n)} = \begin{pmatrix}
			\delta \phi^{(n)}\\[1mm]
			\delta f^{(n)}\\[1mm]
			r^{(n)}\\[1mm]
			a^{(n)}\\[1mm]
			\delta\epsilon_1^{(n)}\\[1mm]
			\vdots\\[1mm]
			\delta\epsilon_{M^{(n)}}^{(n)}
		\end{pmatrix}.
	\end{equation}
	The state evolution is as in the single--pulsar case, except that the gravitational-wave noise processes $\chi_a^{(n)}(t)$ are correlated among pulsars. Their ensemble statistics are
	\begin{equation}
		\langle \chi_a^{(n)}(t) \rangle = 0,
	\end{equation}
	\begin{equation}
		\langle \chi_a^{(n)}(t)\,\chi_a^{(n')}(t') \rangle = \left[\sigma_a^{(n,n')}\right]^2\,\delta(t-t'),
	\end{equation}
	with
	\begin{equation}
		\left[\sigma_a^{(n,n')}\right]^2 = \frac{\langle h^2 \rangle}{6}\,\gamma_a\,\Gamma\bigl[\theta^{(n,n')}\bigr],
	\end{equation}
	and
	\begin{equation}
		\Gamma\bigl[\theta^{(n,n')}\bigr] = \frac{3}{2}x_{nn'}\ln x_{nn'} - \frac{x_{nn'}}{4} + \frac{1}{2} + \frac{1}{2}\delta_{nn'}, \qquad x_{nn'}=\frac{1-\cos\theta^{(n,n')}}{2}\,.
	\end{equation}
	
	\subsection*{1. Discretised State Transition Matrix (\(\boldsymbol{F}\))}
	
	For pulsar $n$, the \((\delta\phi,\delta f)\) block discretises as
	\begin{equation}
		F_p^{(n)} = \begin{pmatrix}
			1 & \dfrac{1-e^{-\gamma_p^{(n)}\Delta t}}{\gamma_p^{(n)}} \\[1mm]
			0 & e^{-\gamma_p^{(n)}\Delta t}
		\end{pmatrix},
	\end{equation}
	and the \((r,a)\) block is
	\begin{equation}
		F_a = \begin{pmatrix}
			1 & \dfrac{1-e^{-\gamma_a\Delta t}}{\gamma_a} \\[1mm]
			0 & e^{-\gamma_a\Delta t}
		\end{pmatrix},
	\end{equation}
	(with the same gravitational–wave damping rate $\gamma_a$ assumed for all pulsars). The timing model parameters evolve by the identity. Hence, the state transition matrix for pulsar $n$ is
	\begin{equation}
		F^{(n)} = \begin{pmatrix}
			F_p^{(n)} & 0 & 0\\[2mm]
			0 & F_a & 0\\[2mm]
			0 & 0 & I_{M^{(n)}}
		\end{pmatrix}.
	\end{equation}
	Stacking the pulsar states, the overall state transition matrix is block diagonal:
	\begin{equation}
		\boldsymbol{F} = \mathrm{diag}\Bigl\{F^{(1)},F^{(2)},\ldots,F^{(N)}\Bigr\}\,.
	\end{equation}
	
	\subsection*{2. Discretised Process Noise Covariance (\(\boldsymbol{Q}\))}
	
	For pulsar $n$, the \((\delta\phi,\delta f)\) block has the discretised covariance
	\begin{equation}
		Q_p^{(n)} = \sigma_p^{(n)2}\begin{pmatrix}
			\displaystyle \frac{\Delta t}{\gamma_p^{(n)2}} - \frac{2\left(1-e^{-\gamma_p^{(n)}\Delta t}\right)}{\gamma_p^{(n)3}} + \frac{1-e^{-2\gamma_p^{(n)}\Delta t}}{2\gamma_p^{(n)3}} & \displaystyle \frac{1-e^{-\gamma_p^{(n)}\Delta t}}{\gamma_p^{(n)2}} - \frac{1-e^{-2\gamma_p^{(n)}\Delta t}}{2\gamma_p^{(n)2}}\\[2mm]
			\displaystyle \frac{1-e^{-\gamma_p^{(n)}\Delta t}}{\gamma_p^{(n)2}} - \frac{1-e^{-2\gamma_p^{(n)}\Delta t}}{2\gamma_p^{(n)2}} & \displaystyle \frac{1-e^{-2\gamma_p^{(n)}\Delta t}}{2\gamma_p^{(n)}}
		\end{pmatrix}.
	\end{equation}
	For the gravitational--wave (or \(a\)) block, note that the continuous noise processes are correlated among different pulsars. Thus, the cross–covariance between pulsars $n$ and $n'$ is discretised as
	\begin{equation}
		Q_a^{(n,n')} = \left[\sigma_a^{(n,n')}\right]^2\begin{pmatrix}
			\displaystyle \frac{\Delta t}{\gamma_a^2} - \frac{2\left(1-e^{-\gamma_a\Delta t}\right)}{\gamma_a^3} + \frac{1-e^{-2\gamma_a\Delta t}}{2\gamma_a^3} & \displaystyle \frac{1-e^{-\gamma_a\Delta t}}{\gamma_a^2} - \frac{1-e^{-2\gamma_a\Delta t}}{2\gamma_a^2}\\[2mm]
			\displaystyle \frac{1-e^{-\gamma_a\Delta t}}{\gamma_a^2} - \frac{1-e^{-2\gamma_a\Delta t}}{2\gamma_a^2} & \displaystyle \frac{1-e^{-2\gamma_a\Delta t}}{2\gamma_a}
		\end{pmatrix}.
	\end{equation}
	For the timing model parameters of pulsar $n$, modeled as a random walk,
	\begin{equation}
		Q_\epsilon^{(n)} = \sigma_\epsilon^2\,\Delta t\,I_{M^{(n)}}.
	\end{equation}
	
	The block coupling pulsars $n$ and $n'$ is then assembled as
	\begin{equation}
		\boldsymbol{Q}^{(n,n')} = \begin{pmatrix}
			\delta_{nn'}\, Q_p^{(n)} & 0 & 0 \\[2mm]
			0 & Q_a^{(n,n')} & 0 \\[2mm]
			0 & 0 & \delta_{nn'}\, Q_\epsilon^{(n)}
		\end{pmatrix},
	\end{equation}
	where $\delta_{nn'}$ is the Kronecker delta (i.e. the spin and timing noise are uncorrelated between different pulsars). Finally, the full process noise covariance is the block matrix
	\begin{equation}
		\boldsymbol{Q} = \begin{pmatrix}
			\boldsymbol{Q}^{(1,1)} & \boldsymbol{Q}^{(1,2)} & \cdots & \boldsymbol{Q}^{(1,N)} \\[1mm]
			\boldsymbol{Q}^{(2,1)} & \boldsymbol{Q}^{(2,2)} & \cdots & \boldsymbol{Q}^{(2,N)} \\[1mm]
			\vdots & \vdots & \ddots & \vdots \\[1mm]
			\boldsymbol{Q}^{(N,1)} & \boldsymbol{Q}^{(N,2)} & \cdots & \boldsymbol{Q}^{(N,N)}
		\end{pmatrix}\,.
	\end{equation}
	
	\subsection*{3. Measurement Matrix (\(\boldsymbol{H}\))}
	
	For pulsar $n$, the measurement equation is
	\begin{equation}
		\delta t^{(n)} = \frac{\delta\phi^{(n)}}{f_0^{(n)}} + \boldsymbol{M}^{(n)}\,\boldsymbol{\delta\epsilon}^{(n)} - r^{(n)}\,,
	\end{equation}
	so that the measurement matrix is
	\begin{equation}
		\boldsymbol{H}^{(n)} = \begin{pmatrix}
			\frac{1}{f_0^{(n)}} & 0 & -1 & 0 & M_1^{(n)} & \cdots & M_{M^{(n)}}^{(n)}
		\end{pmatrix}\,.
	\end{equation}
	In a multi–pulsar analysis the overall measurement vector $\boldsymbol{Y}$ is built by stacking the individual measurements. In practice, since pulsars are generally observed at different times, only a subset of the pulsars contribute at any given time; the full measurement matrix is then formed by selecting the corresponding rows from the block–diagonal matrix
	\begin{equation}
		\boldsymbol{H}_{\mathrm{full}} = \mathrm{diag}\Bigl\{\boldsymbol{H}^{(1)},\boldsymbol{H}^{(2)},\ldots,\boldsymbol{H}^{(N)}\Bigr\}\,.
	\end{equation}
	
	\subsection*{4. Measurement Noise Covariance (\(\boldsymbol{R}\))}
	
	Assuming that the measurement noise for pulsar $n$ is white with variance $\sigma_t^{(n)2}$ and that different pulsars have independent measurement noise, then if at a given time step the set of observed pulsars is $\mathcal{O}\subset\{1,\ldots,N\}$, the measurement noise covariance is
	\begin{equation}
		\boldsymbol{R} = \mathrm{diag}\Bigl\{\sigma_t^{(n)2}\Bigr\}_{n\in\mathcal{O}}\,.
	\end{equation}


	
%\section{additional notes}
%
%It may help the coding if we reorder some things...
%
%\section{Pipeline}
%
%\begin{enumerate}
%	\item Obtain a .par and .tim file for a pulsar
%	\item Pass these files through TEMPO/PINT to obtain timing resiudals $\delta t$, and a design matrix $\mathbf{M}$.
%\end{enumerate}
%
%\section{Dispersion and multiband}
%
%all of the above assumes wideband


\newpage
\appendix
\section{A worked example for the Q-matrix}


It can help intution to "see" what the Q-matrix looks like explicitly. Consider the case where $N=2$. \newline 



	Recall that for each pulsar the state vector is partitioned into three sectors:
	\begin{enumerate}
		\item \textbf{Spin noise sector} (for \(\delta\phi\) and \(\delta f\)): a \(2\times2\) block.
		\item \textbf{Gravitational--wave (redshift/residual) sector} (for \(r\) and \(a\)): a \(2\times2\) block.
		\item \textbf{Timing model (design matrix) sector} (for \(\delta\epsilon\)): a block of size \(M^{(n)}\times M^{(n)}\).
	\end{enumerate}
	
	For pulsar \(n\) (\(n=1,2\)), the individual blocks are defined as follows:
	
	\subsection*{Spin Noise Sector}
	
	For pulsar \(n\), the discretised spin noise covariance is
	\begin{equation}
		Q_p^{(n)} = \sigma_p^{(n)2}
		\begin{pmatrix}
			\displaystyle \frac{\Delta t}{\gamma_p^{(n)2}} - \frac{2\bigl(1-e^{-\gamma_p^{(n)}\Delta t}\bigr)}{\gamma_p^{(n)3}} + \frac{1-e^{-2\gamma_p^{(n)}\Delta t}}{2\gamma_p^{(n)3}} 
			& \displaystyle \frac{1-e^{-\gamma_p^{(n)}\Delta t}}{\gamma_p^{(n)2}} - \frac{1-e^{-2\gamma_p^{(n)}\Delta t}}{2\gamma_p^{(n)2}} \\[2mm]
			\displaystyle \frac{1-e^{-\gamma_p^{(n)}\Delta t}}{\gamma_p^{(n)2}} - \frac{1-e^{-2\gamma_p^{(n)}\Delta t}}{2\gamma_p^{(n)2}} 
			& \displaystyle \frac{1-e^{-2\gamma_p^{(n)}\Delta t}}{2\gamma_p^{(n)}}
		\end{pmatrix}\,.
	\end{equation}
	
	\subsection*{Gravitational--Wave Sector}
	
	For the gravitational--wave noise, the continuous noise processes \(\chi_a^{(n)}(t)\) are correlated among pulsars. Thus, for pulsars \(n\) and \(n'\) the discretised covariance is given by:
	\begin{equation}
		Q_a^{(n,n')} = \left[\sigma_a^{(n,n')}\right]^2
		\begin{pmatrix}
			\displaystyle \frac{\Delta t}{\gamma_a^2} - \frac{2\bigl(1-e^{-\gamma_a\Delta t}\bigr)}{\gamma_a^3} + \frac{1-e^{-2\gamma_a\Delta t}}{2\gamma_a^3} 
			& \displaystyle \frac{1-e^{-\gamma_a\Delta t}}{\gamma_a^2} - \frac{1-e^{-2\gamma_a\Delta t}}{2\gamma_a^2} \\[2mm]
			\displaystyle \frac{1-e^{-\gamma_a\Delta t}}{\gamma_a^2} - \frac{1-e^{-2\gamma_a\Delta t}}{2\gamma_a^2} 
			& \displaystyle \frac{1-e^{-2\gamma_a\Delta t}}{2\gamma_a}
		\end{pmatrix}\,.
	\end{equation}
	
	For the \textbf{autocovariance} (i.e. \(n=n'\)) we denote this as \(Q_a^{(n,n)}\), and for the \textbf{cross-covariance} (i.e. \(n \neq n'\)) we have
	\[
	Q_a^{(1,2)} = Q_a^{(2,1)}.
	\]
	
	\subsection*{Timing Model Sector}
	
	For pulsar \(n\), the timing model parameters are assumed to follow a random walk. Their covariance is given by
	\begin{equation}
		Q_\epsilon^{(n)} = \sigma_\epsilon^2\,\Delta t\,I_{M^{(n)}},
	\end{equation}
	where \(I_{M^{(n)}}\) is the identity matrix of dimension \(M^{(n)}\).
	
	\subsection*{Constructing the Full \(\boldsymbol{Q}\) Matrix for Two Pulsars}
	
	For two pulsars (\(N=2\)), we construct the overall process noise covariance matrix \(\boldsymbol{Q}\) as a \(2\times2\) block matrix:
	\begin{equation}
		\boldsymbol{Q} =
		\begin{pmatrix}
			\boldsymbol{Q}^{(1,1)} & \boldsymbol{Q}^{(1,2)} \\[1mm]
			\boldsymbol{Q}^{(2,1)} & \boldsymbol{Q}^{(2,2)}
		\end{pmatrix}\,.
	\end{equation}
	
	Each block \(\boldsymbol{Q}^{(n,n')}\) is itself block–structured into three sectors (spin noise, gravitational–wave noise, and timing model noise).
	
	\subsubsection*{Diagonal Block for Pulsar 1 (\(\boldsymbol{Q}^{(1,1)}\))}
	
	\begin{equation}
		\boldsymbol{Q}^{(1,1)} =
		\begin{pmatrix}
			Q_p^{(1)} & \boldsymbol{0}_{2\times2} & \boldsymbol{0}_{2\times M^{(1)}} \\[2mm]
			\boldsymbol{0}_{2\times2} & Q_a^{(1,1)} & \boldsymbol{0}_{2\times M^{(1)}} \\[2mm]
			\boldsymbol{0}_{M^{(1)}\times2} & \boldsymbol{0}_{M^{(1)}\times2} & Q_\epsilon^{(1)}
		\end{pmatrix}\,.
	\end{equation}
	
	Here:
	\begin{itemize}
		\item The upper-left \(2\times2\) block is \(Q_p^{(1)}\) (spin noise for pulsar 1).
		\item The middle \(2\times2\) block is \(Q_a^{(1,1)}\) (gravitational–wave autocovariance for pulsar 1).
		\item The lower-right \(M^{(1)}\times M^{(1)}\) block is \(Q_\epsilon^{(1)}\) (timing model noise for pulsar 1).
	\end{itemize}
	
	\subsubsection*{Diagonal Block for Pulsar 2 (\(\boldsymbol{Q}^{(2,2)}\))}
	
	\begin{equation}
		\boldsymbol{Q}^{(2,2)} =
		\begin{pmatrix}
			Q_p^{(2)} & \boldsymbol{0}_{2\times2} & \boldsymbol{0}_{2\times M^{(2)}} \\[2mm]
			\boldsymbol{0}_{2\times2} & Q_a^{(2,2)} & \boldsymbol{0}_{2\times M^{(2)}} \\[2mm]
			\boldsymbol{0}_{M^{(2)}\times2} & \boldsymbol{0}_{M^{(2)}\times2} & Q_\epsilon^{(2)}
		\end{pmatrix}\,.
	\end{equation}
	
	Here:
	\begin{itemize}
		\item The upper-left \(2\times2\) block is \(Q_p^{(2)}\) (spin noise for pulsar 2).
		\item The middle \(2\times2\) block is \(Q_a^{(2,2)}\) (gravitational–wave autocovariance for pulsar 2).
		\item The lower-right \(M^{(2)}\times M^{(2)}\) block is \(Q_\epsilon^{(2)}\) (timing model noise for pulsar 2).
	\end{itemize}
	
	\subsubsection*{Off-Diagonal Block Between Pulsar 1 and Pulsar 2 (\(\boldsymbol{Q}^{(1,2)}\))}
	
	\begin{equation}
		\boldsymbol{Q}^{(1,2)} =
		\begin{pmatrix}
			\boldsymbol{0}_{2\times2} & \boldsymbol{0}_{2\times2} & \boldsymbol{0}_{2\times M^{(2)}} \\[2mm]
			\boldsymbol{0}_{2\times2} & Q_a^{(1,2)} & \boldsymbol{0}_{2\times M^{(2)}} \\[2mm]
			\boldsymbol{0}_{M^{(1)}\times2} & \boldsymbol{0}_{M^{(1)}\times2} & \boldsymbol{0}_{M^{(1)}\times M^{(2)}}
		\end{pmatrix}\,.
	\end{equation}
	
	In this off-diagonal block:
	\begin{itemize}
		\item The spin noise sectors (upper-left \(2\times2\)) are zero because spin noise is uncorrelated between pulsars.
		\item The timing model sectors are zero.
		\item Only the gravitational–wave sector (the middle \(2\times2\) block) is nonzero and given by \(Q_a^{(1,2)}\).
	\end{itemize}
	
	\subsubsection*{Off-Diagonal Block Between Pulsar 2 and Pulsar 1 (\(\boldsymbol{Q}^{(2,1)}\))}
	
	By symmetry (assuming the cross–covariance is symmetric), we have
	\begin{equation}
		\boldsymbol{Q}^{(2,1)} =
		\begin{pmatrix}
			\boldsymbol{0}_{2\times2} & \boldsymbol{0}_{2\times2} & \boldsymbol{0}_{2\times M^{(1)}} \\[2mm]
			\boldsymbol{0}_{2\times2} & Q_a^{(2,1)} & \boldsymbol{0}_{2\times M^{(1)}} \\[2mm]
			\boldsymbol{0}_{M^{(2)}\times2} & \boldsymbol{0}_{M^{(2)}\times2} & \boldsymbol{0}_{M^{(2)}\times M^{(1)}}
		\end{pmatrix}\,,
	\end{equation}
	with \(Q_a^{(2,1)} = Q_a^{(1,2)}\).
	
	\subsection*{Full \(\boldsymbol{Q}\) Matrix for \(N=2\)}
	
	Assembling the blocks, the full \( \boldsymbol{Q} \) matrix is given by
	\begin{equation}
		\boldsymbol{Q} =
		\begin{pmatrix}
			\begin{array}{c|c}
				\begin{array}{ccc}
					Q_p^{(1)} & \boldsymbol{0}_{2\times2} & \boldsymbol{0}_{2\times M^{(1)}}
					\\[2mm]
					\boldsymbol{0}_{2\times2} & Q_a^{(1,1)} & \boldsymbol{0}_{2\times M^{(1)}}
					\\[2mm]
					\boldsymbol{0}_{M^{(1)}\times2} & \boldsymbol{0}_{M^{(1)}\times2} & Q_\epsilon^{(1)}
				\end{array}
				& 
				\begin{array}{ccc}
					\boldsymbol{0}_{2\times2} & \boldsymbol{0}_{2\times2} & \boldsymbol{0}_{2\times M^{(2)}
					}\\[2mm]
					\boldsymbol{0}_{2\times2} & Q_a^{(1,2)} & \boldsymbol{0}_{2\times M^{(2)}
					}\\[2mm]
					\boldsymbol{0}_{M^{(1)}\times2} & \boldsymbol{0}_{M^{(1)}\times2} & \boldsymbol{0}_{M^{(1)}\times M^{(2)}
					}
				\end{array}
				\\[4mm]
				\hline \\[-3mm]
				\begin{array}{ccc}
					\boldsymbol{0}_{2\times2} & \boldsymbol{0}_{2\times2} & \boldsymbol{0}_{2\times M^{(1)}
					}\\[2mm]
					\boldsymbol{0}_{2\times2} & Q_a^{(2,1)} & \boldsymbol{0}_{2\times M^{(1)}
					}\\[2mm]
					\boldsymbol{0}_{M^{(2)}\times2} & \boldsymbol{0}_{M^{(2)}\times2} & \boldsymbol{0}_{M^{(2)}\times M^{(1)}
					}
				\end{array}
				&
				\begin{array}{ccc}
					Q_p^{(2)} & \boldsymbol{0}_{2\times2} & \boldsymbol{0}_{2\times M^{(2)}
					}\\[2mm]
					\boldsymbol{0}_{2\times2} & Q_a^{(2,2)} & \boldsymbol{0}_{2\times M^{(2)}
					}\\[2mm]
					\boldsymbol{0}_{M^{(2)}\times2} & \boldsymbol{0}_{M^{(2)}\times2} & Q_\epsilon^{(2)}
				\end{array}
			\end{array}
		\end{pmatrix}\,.
	\end{equation}
	
	\textbf{Component Summary:}
	\begin{itemize}
		\item \(\boldsymbol{Q}^{(1,1)}\) (Pulsar 1):
		\begin{itemize}
			\item \textbf{Spin noise:} \(Q_p^{(1)}\) (upper-left \(2\times2\))
			\item \textbf{Gravitational--wave noise:} \(Q_a^{(1,1)}\) (middle \(2\times2\))
			\item \textbf{Timing model noise:} \(Q_\epsilon^{(1)}\) (lower-right \(M^{(1)}\times M^{(1)}\))
		\end{itemize}
		\item \(\boldsymbol{Q}^{(2,2)}\) (Pulsar 2):
		\begin{itemize}
			\item \textbf{Spin noise:} \(Q_p^{(2)}\)
			\item \textbf{Gravitational--wave noise:} \(Q_a^{(2,2)}\)
			\item \textbf{Timing model noise:} \(Q_\epsilon^{(2)}\)
		\end{itemize}
		\item \(\boldsymbol{Q}^{(1,2)}\) (Between Pulsar 1 and Pulsar 2):
		\begin{itemize}
			\item \textbf{Spin noise:} Zero (\(2\times2\) zero matrix)
			\item \textbf{Gravitational--wave noise:} \(Q_a^{(1,2)}\) (middle \(2\times2\))
			\item \textbf{Timing model noise:} Zero
		\end{itemize}
		\item \(\boldsymbol{Q}^{(2,1)}\) (Between Pulsar 2 and Pulsar 1):
		\begin{itemize}
			\item \textbf{Spin noise:} Zero
			\item \textbf{Gravitational--wave noise:} \(Q_a^{(2,1)}\) (which equals \(Q_a^{(1,2)}\))
			\item \textbf{Timing model noise:} Zero
		\end{itemize}
	\end{itemize}
	
	This explicit layout shows how the overall \( \boldsymbol{Q} \) matrix incorporates:
	\begin{itemize}
		\item The individual (uncorrelated) spin noise and timing model noise (appearing in the diagonal blocks).
		\item The cross–correlations due to the gravitational--wave background (appearing in the off-diagonal blocks).
	\end{itemize}
	



\end{document}
